\documentclass[11pt,a4paper]{article}
\usepackage[utf8]{inputenc}
\usepackage[T1]{fontenc}
\usepackage[lithuanian]{babel}
\usepackage[margin=1.5 cm]{geometry}
\usepackage{longtable}
\usepackage{array}
\usepackage[table]{xcolor}
\usepackage{fancyhdr}

\pagestyle{fancy}
\fancyhf{}
\fancyfoot[C]{\thepage}
\renewcommand{\headrulewidth}{0pt}
\fancypagestyle{plain}{%
  \fancyhf{}
  \fancyfoot[C]{\thepage}
  \renewcommand{\headrulewidth}{0pt}
}

\begin{document}

\pagenumbering{arabic}

\begin{flushright}
PATVIRTINTA \\
Vilniaus licėjaus direktoriaus \\
2023 m. sausio 31 d. įsakymu V1-21
\end{flushright}

\begin{center}
\textbf{VILNIAUS LICĖJUS} \\
\textbf{2026–2027 m. m. ILGALAIKIS PLANAS}
\end{center}

\noindent
\begin{tabular}{@{}l@{\hspace{1.5em}}l@{}}
\textbf{DALYKAS} & Matematika \\
\textbf{KLASĖ} & I klasė \\
\textbf{PAMOKŲ SKAIČIUS} & 148 (4 pamokos per savaitę)
\end{tabular}

\vspace{0.5cm}

\noindent \textbf{TIKSLAS} – sudaryti galimybę kiekvienam mokiniui, mokantis matematikos, ugdytis matematinį ir statistinį raštingumą, kuris šiame dokumente suprantamas kaip įgytas gebėjimas matematiškai samprotauti ir taikyti įgytas kompetencijas, sprendžiant įvairias realias, aktualias ir mokiniams suprantamas problemas.

\vspace{0.3cm}
\noindent \textbf{UŽDAVINIAI} \\
Siekdami tikslo mokiniai:
\begin{itemize}
    \item tinkamai ir tikslingai vartoja matematinius faktus; sklandžiai atlieka matematines procedūras; įgytas žinias sieja tarpusavyje, sistemina, struktūruoja; įžvelgia matematikos ryšius su kitais dalykais;
    \item įvairiuose kontekstuose taiko indukcinį ir dedukcinį, kiekybinį ir statistinį samprotavimą; remiasi žiniomis, logika ir patikimais argumentais, formuluodami, analizuodami, įrodinėdami teiginius, spręsdami uždavinius, darydami išvadas ar vertinimus;
    \item bendradarbiaudami su kitais, nagrinėja įvairiomis formomis pateiktus matematinius pranešimus, dalyvauja diskusijose apie komunikavimo tikslą, adresatą, pranešimu perteikiamų minčių tikslumą, logiškumą, pagrįstumą, išsamumą, glaustumą;
    \item yra nusiteikę ir įdeda pastangų matematikos mokymosi kliūtims įveikti; tikslingai planuoja ir organizuoja mokymosi veiklą; siekdami mokytis matematikos ir ją pažinti, turi žinių, gebėjimų ir polinkį naudotis skaitmeninėmis technologijomis;
    \item įgytas matematines kompetencijas ir supratimą apie bendrą problemų sprendimo procesą kūrybiškai pritaiko įvairiuose realiuose, aktualiuose ir mokiniams suprantamuose kontekstuose; reflektuoja savo žinias, gebėjimus, samprotavimo veiklą ir jos rezultatus.
\end{itemize}

\vspace{0.3cm}
\noindent \textbf{MOKYMO IR MOKYMOSI PRIEMONĖS}
\begin{itemize}
    \item Matematika 9 klasei. 1, 2 vadovėlio dalys (serija „Horizontai“). Jūratė Gedminienė, Irena Šukienė, Daiva Riukienė, Jolanta Jačiauskaitė, Ingrida Brazauskienė. Kaunas: „Šviesa“, 2023. (prieiga „Eduka klasėje“).
    \item Matematika. Vadovėlis 9 klasei. 1, 2 vadovėlio dalys (serija „TEMPUS”). Jolanta Knyvienė, Antanas Apynis, Ingrida Brazauskienė, Ramunė Dranseikienė, Nomeda Zdanienė. Kaunas: „Šviesa“, 2020.
    \item Matematika. Vadovėlis gimnazijų I klasei. 1, 2 vadovėlio dalys. Alvyda Ambraškienė, Rita Belkevičienė, Rita Grigelienė, Birutė Vosylienė, Milda Vosylienė. Kaunas: „Šviesa”, 2008.
\end{itemize}

\vspace{0.3cm}
\noindent \textbf{VERTINIMAS} \\
Mokinių pasiekimai vertinami vadovaujantis Vilniaus licėjaus mokinių pasiekimų ir pažangos vertinimo tvarkos aprašu, patvirtintu direktoriaus 2024-06-27 įsakymu Nr. V1-112. \\
Atsižvelgiant į pamokos uždavinius, nuolat taikomas formuojamasis vertinimas. Už aktyvumą pamokos metu, atliktas namų darbų, klasės darbų, papildomas užduotis mokiniams skiriami kaupiamojo įvertinimo balai. Surinkus 10 kaupiamųjų įvertinimo balų, į dienyną rašomas 10 balų įvertinimas. Kiekvieno skyriaus pabaigoje taikomas apibendrinamasis įvertinimas (kontrolinis darbas). Apibendrinamasis vertinimas (savarankiškas darbas, apklausa) taip pat yra taikomas baigiant svarbią, sudėtingą temą. Diagnostinės užduotys parengiamos atsižvelgiant į matematikos bendrojoje programoje numatytus pasiekimus, pasiekimų lygius, žinių ir gebėjimų santykį.

\vspace{0.5cm}

\renewcommand{\arraystretch}{1.5}
\begin{longtable}{|>{\raggedright\arraybackslash}p{3.5cm}|>{\centering\arraybackslash}p{0.8cm}|p{2.3cm}|p{6.5cm}|p{2cm}|}
\hline
\multicolumn{1}{|c|}{\textbf{Pamokos tema}} & \multicolumn{1}{c|}{\textbf{Val.}} & \multicolumn{1}{c|}{\textbf{Vertinimas}} & \multicolumn{1}{c|}{\textbf{Kompetencijos ir veiklos}} & \parbox[c][2.4em][c]{2cm}{\centering \rule{0pt}{1.1em}\textbf{Plano}\\[-0.3em]\textbf{korekcija}} \\ \hline
\endfirsthead

\hline
\multicolumn{1}{|c|}{\textbf{Pamokos tema}} & \multicolumn{1}{c|}{\textbf{Val.}} & \multicolumn{1}{c|}{\textbf{Vertinimas}} & \multicolumn{1}{c|}{\textbf{Kompetencijos ir veiklos}} & \parbox[c][2.4em][c]{2cm}{\centering \rule{0pt}{1.1em}\textbf{Plano}\\[-0.3em]\textbf{korekcija}} \\ \hline
\endhead

\hline
\endfoot

\hline
\endlastfoot

\rowcolor{gray!20}
\multicolumn{4}{|l|}{\textbf{1 Tema. Daugianariai}} & \textbf{Val.sk. 10} \\ \hline
Algebrinių reiškinių tipai & 1 & & Socialinė, emocinė ir sveikos gyvensenos kompetencija (mokomasi įsivertinti ir planuoti tolimesnį mokymąsi). Pažinimo kompetencija. & \\ \hline
Laipsnis su natūraliuoju rodikliu ir jo savybės & 1 & & & \\ \hline
Laipsnis su sveikuoju rodikliu ir jo savybės & 1 & & & \\ \hline
Vienanariai ir veiksmai su jais & 1 & & & \\ \hline
Daugianaris & 1 & & & \\ \hline
Daugianarių sudėtis, atimtis, daugyba & 1 & & & \\ \hline
Greitosios daugybos formulės & 1 & & & \\ \hline
Pilnojo kvadrato išskyrimas & 1 & & & \\ \hline
Kvadratinio daugianario ekstremumai & 1 & & & \\ \hline
Kontrolinis darbas & 1 & Kontrolinis darbas & & \\ \hline

\rowcolor{gray!20}
\multicolumn{4}{|l|}{\textbf{2 Tema. Kvadratinės lygtys}} & \textbf{Val.sk. 7} \\ \hline
Nepilnoji kvadratinė lygtis & 1 & & & \\ \hline
Pilnoji kvadratinė lygtis & 1 & & & \\ \hline
Vijeto teorema & 1 & & & \\ \hline
Kvadratinio trinario skaidymas dauginamaisiais & 1 & & & \\ \hline
Bikvadratinė lygtis & 1 & & & \\ \hline
Kvadratinė lygtis įvedant keitinius & 1 & & & \\ \hline
Nestandartiniai uždaviniai su diskriminantu ir Vijeto teorema & 1 & & & \\ \hline
Kontrolinis darbas & 1 & Kontrolinis darbas & & \\ \hline

\rowcolor{gray!20}
\multicolumn{4}{|l|}{\textbf{3 Tema. Lygčių sistemos}} & \textbf{Val.sk. 4} \\ \hline
Tiesinė lygčių sistema ir jos sprendimas & 2 & & & \\ \hline
Netiesinė lygčių sistema ir jos sprendimas & 1 & & & \\ \hline
Tekstinių uždavinių sprendimas sudarant lygčių sistemą & 3 & & & \\ \hline
Simetriniai daugianariai ir lygčių sistemos & 2 & & & \\ \hline\
Kontrolinis darbas & 1 & Kontrolinis darbas & & \\ \hline
\rowcolor{gray!20}
\multicolumn{4}{|l|}{\textbf{4 Tema. Geometrija}} & \textbf{Val.sk. 12} \\ \hline
Tiesinių lygčių sistemų sprendimas keitimo būdu & 2 & & Pažinimo kompetencija (apibrėžiamos sąvokos: lygtis su dviem nežinomaisiais, lygties su dviem nežinomaisiais sprendinys, mokomasi spręsti tiesinių lygčių sistemos, mokomasi spręsti sistemas, kuriose viena lygtis tiesinė, o kita - kvadratinė). Kūrybiškumo kompetencija, pilietiškumo kompetencija (nagrinėjamos įvairios realaus pasaulio situacijos, kurios gali būti modeliuojamos lygčių sistemomis). & \\ \hline
Tiesinių lygčių sistemų sprendimas sudėties būdu & 1 & & & \\ \hline
Tiesinių lygčių sudarymas & 4 & Rašto darbas & & \\ \hline
Lygčių sistemos, kurių tik viena lygtis tiesinė & 1 & & & \\ \hline
Tekstiniai uždaviniai & 3 & & & \\ \hline
Kontrolinis darbas & 1 & Kontrolinis darbas & & \\ \hline

\rowcolor{gray!20}
\multicolumn{4}{|l|}{\textbf{5 Tema. Trikampio geometrija}} & \textbf{Val.sk. 7} \\ \hline
Plotų metodas & 1 & & & \\ \hline
Talio teorema & 1 & & & \\ \hline
Atvirkštinė Talio teorema & 1 & & & \\ \hline
Pusiaukraštinių susikirtimo taškas ir jo savybės & 1 & & & \\ \hline
Pusiaukampinių susikirtimo taškas ir jo savybės & 1 & & & \\ \hline
Aukštinių susikirtimo taškas & 1 & & & \\ \hline
Brėžimo uždaviniai & 1 & & & \\ \hline

\rowcolor{gray!20}
\multicolumn{4}{|l|}{\textbf{6 Tema. Panašieji trikampiai}} & \textbf{Val.sk. 17} \\ \hline
Trikampių kartojimas & 2 & & Pažinimo kompetencija (apibrėžiama, kokios figūros vadinamos panašiosiomis, taikomi trikampių panašumo požymiai). & \\ \hline
Keturkampių kartojimas & 2 & & & \\ \hline
Proporcingosios atkarpos & 1 & & & \\ \hline
Panašiųjų trikampių savybės & 3 & & & \\ \hline
Trikampių panašumo požymiai & 4 & Rašto darbas & & \\ \hline
Trikampio vidurinė linija & 1 & & & \\ \hline
Trapecijos vidurinė linija & 1 & & & \\ \hline
Trikampių panašumo požymių taikymas & 2 & & & \\ \hline
Kontrolinis darbas & 1 & Kontrolinis darbas & & \\ \hline

\rowcolor{gray!20}
\multicolumn{4}{|l|}{\textbf{7 Tema. Apskritimo geometrija}} & \textbf{Val.sk. 4} \\ \hline
Centriniai ir įbrėžtiniai kampai ir jų sąryšis & 1 & & & \\ \hline
Liestinė ir kirstinė & 1 & & & \\ \hline
Skritulio išpjova ir nuopjova & 1 & & & \\ \hline
Apskritimai ir laikrodžiai & 1 & & & \\ \hline

\rowcolor{gray!20}
\multicolumn{4}{|l|}{\textbf{8 Tema. Funkcijos}} & \textbf{Val.sk. 7} \\ \hline
Funkcija ir jos reiškimo būdai & 2 & & Pažinimo kompetencija (apibrėžiamos funkcijos, funkcijos argumento, funkcijos reikšmės, funkcijos apibrėžimo srities, funkcijos reikšmių srities sąvokos. Komunikavimo kompetencija (mokomasi funkciją apibūdinti žodžiais, lentele, grafiku formule). Skaitmeninė kompetencija (mokomasi apibūdinti funkciją naudojantis skaitmeninėmis priemonėmis). & \\ \hline
Funkcijos savybės ir jos grafikas & 2 & & & \\ \hline
Funkcijos savybės ir jos formulė & 3 & Rašto darbas & & \\ \hline

\rowcolor{gray!20}
\multicolumn{4}{|l|}{\textbf{9 Tema. Tiesinė funkcija}} & \textbf{Val.sk. 9} \\ \hline
Tiesinė funkcija ir jos grafikas & 2 & & Pažinimo kompetencija (apibrėžiama tiesinė funkcija, funkcijos krypties koeficientas). Skaitmeninė kompetencija (naudojantis skaitmeninėmis priemonėmis, nagrinėjama, kaip keičiasi funkcijos grafikas taikant transformacijas). Pilietiškumo kompetencija (sprendžiami uždaviniai, kuriuose įvairios realaus konteksto situacijos yra modeliuojamos funkcijomis). & \\ \hline
Tiesinės funkcijos savybės & 2 & & & \\ \hline
Tiesinės funkcijos formulės sudarymas & 2 & & & \\ \hline
Tiesinės funkcijos taikymas gyvenimiškose situacijose & 2 & & & \\ \hline
Kontrolinis darbas & 1 & Kontrolinis darbas & & \\ \hline

\rowcolor{gray!20}
\multicolumn{4}{|l|}{\textbf{10 Tema. Kvadratinės funkcijos}} & \textbf{Val.sk. 15} \\ \hline
Kvadratinė funkcija $f(x) = ax^2$ & 2 & & Pažinimo kompetencija (įvedama kvadratinės funkcijos sąvoka). Skaitmeninė kompetencija (naudojantis skaitmeninėmis priemonėmis, nagrinėjama, kaip keičiasi funkcijos grafikas taikant transformacijas). Pilietiškumo kompetencija (sprendžiami uždaviniai, kuriuose įvairios realaus konteksto situacijos yra modeliuojamos funkcijomis). & \\ \hline
Kvadratinė funkcija $f(x) = ax^2 + c$ & 3 & Rašto darbas & & \\ \hline
Kvadratinės funkcijos $f(x) = a(x - m)^2 + n$ ir $f(x) = a(x - x_1)(x - x_2)$ & 2 & & & \\ \hline
Kvadratinė funkcija $f(x) = ax^2 + bx + c$ & 3 & & & \\ \hline
Kvadratinių funkcijų taikymas gyvenimiškose situacijose & 2 & & & \\ \hline
Grafinis lygčių sistemų sprendimas & 2 & & & \\ \hline
Kontrolinis darbas & 1 & Kontrolinis darbas & & \\ \hline

\rowcolor{gray!20}
\multicolumn{4}{|l|}{\textbf{11 Tema. Skaičių sekos}} & \textbf{Val.sk. 8} \\ \hline
Skaičių seka & 2 & & Pažinimo kompetencija (skaičių seka apibrėžiama kaip funkcija, paprastais atvejais mokomasi sekas aprašyti n-ojo nario formule ir rekurentiniu būdu). & \\ \hline
Skaičių sekų užrašymas formulėmis. Matematinės indukcijos metodas. & 3 & & & \\ \hline
Skaičių seka -- natūraliojo argumento funkcija & 2 & & & \\ \hline
Kontrolinis darbas & 1 & Kontrolinis darbas & & \\ \hline

\rowcolor{gray!20}
\multicolumn{4}{|l|}{\textbf{12 Tema. Statistika}} & \textbf{Val.sk. 15} \\ \hline
Tikimybės & 2 & & Pažinimo kompetencija (apibrėžiamos populiacijos, imties, imties dydžio, kvartilio sąvokos, nagrinėjamos sklaidos diagramos, mokomasi surasti pirmąjį, antrąjį ir trečiąjį kvartilius, interpretuoti duomenis, kai yra išskirčių). Pilietiškumo kompetencija (nagrinėjant konkrečias situacijas, aptariami imčių sudarymo ir gautų išvadų apie jas pagrįstumo klausimai). Skaitmeninė kompetencija (duomenys analizuojami ir sisteminami naudojant skaitmenines priemones). & \\ \hline
Imtis. Imties atranka & 2 & & & \\ \hline
Imčių sudarymo būdai & 2 & & & \\ \hline
Sklaidos diagrama & 1 & & & \\ \hline
Sklaidos diagramoje pavaizduotos tiesės lygtis & 2 & & & \\ \hline
Koreliacija & 2 & & & \\ \hline
Projektinis darbas & 4 & Projektinis darbas & & \\ \hline

\rowcolor{gray!20}
\multicolumn{4}{|l|}{\textbf{Rezervinės pamokos}} & \textbf{4} \\ \hline

\end{longtable}

\vspace{1cm}
\noindent
\begin{minipage}[t]{0.5\textwidth}
Parengė matematikos mokytojas \\
Andrius Berniukevičius \\
\\
APTARTA \\
xxxxxxxx metodinėje grupėje \\
Protokolo Nr. 1 \\
2024-08- \\
\rule{4cm}{0.4pt} \\
(parašas) \\
Metodinės grupės pirmininkas \\
Inga Vaičiulė
\end{minipage}%
\begin{minipage}[t]{0.5\textwidth}
SUDERINTA \\
\\
\rule{4cm}{0.4pt} \\
(parašas)
\end{minipage}

\end{document}
